% !TEX encoding = UTF-8 Unicode

\documentclass[a4paper,12pt]{report}
\usepackage[natbibapa,nosectionbib,tocbib,numberedbib]{apacite}
\AtBeginDocument{\renewcommand{\bibname}{Literature}}


\usepackage[utf8x]{inputenc}
\usepackage{graphicx}
\usepackage{enumerate}
\usepackage{url}

\usepackage{booktabs}
% \usepackage{graphicx}

\usepackage[colorinlistoftodos]{todonotes}

\usepackage{pifont}

\usepackage{lmodern}
\usepackage{listings}
\lstset{
basicstyle=\scriptsize\ttfamily,
columns=flexible,
breaklines=true,
numbers=left,
%stepsize=1,
numberstyle=\tiny,
backgroundcolor=\color[rgb]{0.85,0.90,1}
}

\usepackage{caption} 
\captionsetup[table]{skip=10pt}

\let\oldquote\quote
\let\endoldquote\endquote
\renewenvironment{quote}{\footnotesize\oldquote}{\endoldquote}



\title{Big Data and Automated Content Analysis\\ (6 ECTS)\\~\\Cursusdossier}
\author{
	dr. Anne Kroon\\
	Graduate School of Communication\\
	University of Amsterdam\\
	a.c.kroon@uva.nl
}

\date{Academic Year 2025/26\\Semester 1, block  2\ }


\begin{document}
\maketitle

\tableofcontents


\chapter{Short description of the course}

``Big data'' refers to data that are more voluminous, but often also more varied, unstructured, and dynamic than data traditionally used in communication research. In Communication Science and the Social Sciences more broadly, this includes research using internet-based data sources such as social-media content, digital news archives, platform data, and public comments. This area of research is commonly described as \emph{Computational Social Science} \citep{Lazer2009} or, when focused on communication processes and content, \emph{Computational Communication Science} \citep{Shah2015}.

This course introduces the concepts, opportunities, and challenges of computational research with digital communication data. Computational approaches can make it possible to analyse large or complex collections of texts and other digital traces; they complement rather than replace established communication-science methods such as theory-driven research design, manual content analysis, and statistical inference. Throughout the course, we emphasise that automated outputs require substantive interpretation and empirical validation.

Participants learn practical strategies for acquiring, organising, preprocessing, analysing, and documenting digital data in communication contexts. The course focuses on (a) data acquisition from files and APIs, data wrangling, exploratory analysis, and reproducible workflows; and (b) automated content analysis, including text preprocessing, dictionary and Bag-of-Words approaches, unsupervised text analysis, embeddings, topic modelling, and supervised machine-learning classification. Students learn to inspect and evaluate automated outputs, compare them with appropriate benchmarks, and reflect on validity, ethics, reproducibility, and the responsible use of AI tools.

The course uses Python as its main programming environment. Students are expected to be willing to learn to write, read, adapt, and explain Python code. Previous programming experience is helpful but not required. Students without prior experience are encouraged to complete one of the many free online introductions to Python before the course begins.


\chapter{Exit qualifications}
The course contributes to the following three exit qualifications of the Research Master's programme in Communication Science:

\textit{Expertise in empirical research}

\begin{enumerate}
    \setcounter{enumi}{2}
    \item Knowledge and Understanding: Have in-depth knowledge and a thorough understanding of advanced research designs and methods.
    \item Skills and abilities: Are able, independently and on their own, to set up, conduct, report, and interpret advanced academic research.
\end{enumerate}

\textit{Academic abilities and attitudes}

\begin{enumerate}
    \setcounter{enumi}{5}
    \item Academic attitudes: Accept that scientific knowledge is always `work in progress' and that arguments must be considered and conclusions drawn on the basis of empirical results and valid criticism.
\end{enumerate}

The course contributes to these exit qualifications at the level of \emph{fundamental} techniques of Big Data and Automated Content Analysis. In particular, students learn to justify methodological choices, acknowledge limitations and uncertainties in their own research, and respond constructively to critical questions. These academic attitudes are assessed through the individual project pitch and the final individual project.

The exit qualifications are elaborated in the following specifications:

\begin{enumerate}
    \setcounter{enumi}{2}
    \item Knowledge and Understanding: Have in-depth knowledge and a thorough understanding of advanced research designs and methods.
    \begin{enumerate}
        \item[3.1.] Have in-depth knowledge and a thorough understanding of advanced research designs and methods, including their value and limitations.
        \item[3.2.] Have in-depth knowledge and a thorough understanding of advanced techniques for data analysis.
    \end{enumerate}

    \item Skills and abilities: Are able, independently and on their own, to set up, conduct, report, and interpret advanced academic research.
    \begin{enumerate}
        \item[4.1.] Are able to formulate research questions and hypotheses for advanced empirical studies.
        \item[4.2.] Are able to develop a research plan, choose appropriate and suitable research designs and methods for advanced empirical studies, and justify the underlying choices.
        \item[4.3.] Are able to assess the validity and reliability of advanced empirical research, and to judge the scientific and professional value of findings from advanced empirical research.
        \item[4.4.] Are able to apply advanced empirical research methods.
    \end{enumerate}

    \item Academic attitudes.
    \begin{enumerate}
        \item[6.1.] Regularly assess their own assumptions, strengths, and weaknesses critically.
        \item[6.2.] Accept that scientific knowledge is always `work in progress' and that something regarded as `true' may be proven to be false, and vice versa.
        \item[6.3.] Are keen to acquire new knowledge, skills, and abilities.
        \item[6.4.] Are willing to share and discuss arguments, results, and conclusions, including submitting one's own work to peer review.
        \item[6.5.] Are convinced that academic debates should not be conducted on the basis of rhetorical qualities but that arguments must be considered and conclusions drawn on the basis of empirical results and valid criticism.
    \end{enumerate}
\end{enumerate}


\chapter{Testable objectives}
\input{testableobjectives}

\chapter{Planning of testing and teaching}

The seminar consists of 12 meetings, two per week. Each week, in the first meeting, the instructor will give short lectures on the key aspects of the week, followed by seminar-style discussions. Theoretical considerations regarding Big Data and Automated Content Analysis are discussed, and techniques for analyzing Big Data are presented. We also discuss examples from the literature, in which these techniques are applied.

The second meetings each week are practicum-meetings, in which the students will apply what the techniques they have learned to their own or provided data sets. Here, they can also deepen their understanding of software tools, prepare their individual projects and get hands-on help. While there are in-class assignments as well as occasional assignments for at home (e.g., completing an online-tutorial to prepare for class), these are not graded.

To complete the course, next to active participation, the students have to successfully complete two summative graded assignments: one mid-term take-home exam and an individual project, in which they derive an empirical question from a theoretical starting point, and then conduct an Automated Content Analysis to answer the question. See Chapter 7 for details.

\chapter{Literature}

\input{literature}

\chapter{Specific course timetable}

\input{schedule}

\chapter{Testing}
An overview of the testing is given in the Test Matrix, displayed in Table \ref{testmatrix}.

\begingroup

\setlength{\tabcolsep}{10pt} % Default value: 6pt
\renewcommand{\arraystretch}{1.5} % Default value: 1

\begin{table}[hbt!]
	\centering
	\caption{Test Matrix}
	\label{testmatrix}
	\resizebox{\columnwidth}{!}{%
	\begin{tabular}{@{}lccc@{}}
	\toprule
	 &
	  \textbf{\begin{tabular}[c]{@{}c@{}}In-class assignments,\\  reviewing work of \\ fellow students,\\  active participation\end{tabular}} &
	  \textbf{\begin{tabular}[c]{@{}c@{}}Mid-term \\ take home exam\end{tabular}} &
	  \textbf{\begin{tabular}[c]{@{}c@{}}Final individual \\ project\end{tabular}} \\ \midrule
	 &
	  \textit{precondition} &
	  \textit{30\% of final grade} &
	  \textit{70\% of final grade} \\
	\begin{tabular}[c]{@{}l@{}}A. Students can explain fundamental research designs and commonly em-\\ ployed methods in existing research articles on Big Data and automated\\ content analysis\end{tabular} &
	  X &
	  X &
	   \\
	\begin{tabular}[c]{@{}l@{}}B. Students can on their own and in own words critically discuss the pros and\\ cons of fundamental research designs and methods employed in existing\\ research articles on Big Data and automated content analysis; they can,\\ based on this, give a critical evaluation of the methods and, where relevant,\\ give advice to improve the study in question.\end{tabular} &
	  X &
	  X &
	   \\
	\begin{tabular}[c]{@{}l@{}}C. Students are able to identify some basic techniques from the field of com-\\ puter science and computer linguistics that are applicable to research in\\ communication science; they can explain the principle of some traditional\\ approaches to text analysis, namely simple rule-based techniques and ba-\\ sic methods of unsupervised and supervised machine learning.\end{tabular} &
	  X &
	  X &
	  X \\
	\begin{tabular}[c]{@{}l@{}}D. Students can on their own formulate a research question and hypotheses\\ for own empirical research in the domain of Big Data.\end{tabular} &
	   &
	   &
	  X \\
	\begin{tabular}[c]{@{}l@{}}E. Students can on their own chose, execute and report on fundamental re-\\ search methods in the domain of Big Data and automatic content analysis.\end{tabular} &
	   &
	   &
	  X \\
	\begin{tabular}[c]{@{}l@{}}F. Students know how to collect data with APIs or read in existing data\\ files; they know how to analyze these data with fundamental automated\\ techniques and to this end, they have basic knowledge of the programming\\ language Python and know how to use fundamental Python-modules for\\ communication science research.\end{tabular} &
	   &
	   &
	  X \\
	\begin{tabular}[c]{@{}l@{}}G Students can critically discuss strong and weak points of their own research\\ using fundamental techniques from the field of Big Data and Automated\\ Content Analysis, and suggest improvements.\end{tabular} &
	   &
	   &
	  X \\
	\begin{tabular}[c]{@{}l@{}}H Students participate actively: reading the literature carefully and on time,\\ completing assignments carefully and on time, active participation in dis-\\ cussions, and giving feedback on the work of fellow students give evidence\\ of this.\end{tabular} &
	  X &
	   &
	   \\ \bottomrule
	\end{tabular}%
	}
	\end{table}

\endgroup

\section*{Grading}
\input{grading}

\chapter{Lecturers’ team, including division of responsibilities}

The course is taught and coordinated by Dr. Anne Kroon.  
She is responsible for all aspects of the course, including teaching, coordination, and the grading of assignments and exams.


\chapter{Calculation of students' study load (in hours)}
\begin{itemize}

\item Elective total: 6 ECTS = 168 hours
\item Reading: 
\begin{itemize}
\item 10 chapters, 5 articles, average 18.5 pages: around 280 pages. 6 pages per hour, thus 46.6 hours for the literature
\item Reading and doing tutorials: 40 hours for reading tutorials to acquire skills.
\item Reading/preparation total: 86.6 hours.
\end{itemize}
\item Presence: \\12*2 hours: 24 hours.
\item Mid-term take-home exam, including preparation: 18 hours
\item Final individual project, including data collection, analysis, write up: 39.4 hours
\end{itemize}

Total: 168 hours


\chapter{Calculation of lecturers' teaching load (in hours)}


\begin{itemize}
\item Presence: 24 hours (= 12 * 2 hours)
\item Preparation of weekly lectures, 7 * 4 hours: 28 hours
\item Preparation of weekly tutorials, 7 * 4 hours: 28 hours
\item Assisting students with setting up Python installation, individual help: 40 hours
\item Feedback and grading take-home exams: 25*60 minutes: 25 hours
\item Feedback and grading final projects, including feedback on proposal and individual counseling: 25* 60 min: 25 hours 
\end{itemize}

Total: 131 hours


\chapter{History of the course}
In response to the feedback by the test review committee in 2019, the following changes were applied:
\begin{itemize}
\item Empasized in the section ``Testable Objectives'' that -- in contrast to the 12 ECTS course -- students should be able to demonstrate, reflect, and apply knowledge about \emph{fundamental} --rather than \emph{advanced} techniques of automated content analysis. Likewise, testing will be based on the extent that students demonstrate knowledge and skills on the level of \emph{fundamental} techniques. This sets the 6ec course apart from Big Data and Automated Content Analysis Part I and Part II, where students take a deeper dive into more advanced and recent techniques for automated content analysis.
\end{itemize}

In 2022, the course received a thorough update due to the publication of the new textbook \citep{cssbook}.

In 2023, the course was slightly updated to incorporate recent developments in machine learning. Most elements remain the same--as the 6ECTS course focuses on fundamental techniques of automated content analysis. Yet, in week 7, we will discuss state of the art techniques for topic modelling. In addition, an outlook is provided on techniques that are the current becoming dominant in the field of machine learning. 

In 2024, the course was slightly updated to make the course more AI proof (using hours from a TLC grant). The two main changes are: 1) The take-home exam will be more focused on debugging and explaining existing code over writing code and the testing will include a short video-walkthrough in addition to the written deliverable. This is aimed at ensuring some level of active engagement with the content by the students without having to use up addional in-class time. 2) One of the last sessions is being used to have students do a small pitch about their final project in class. This has two reasons: First, the students see what their classmates are working on and can exchange ideas and second they can give and receive feedback on the final projects with enough time to still make changes and adjustments.

In 2025, we decided to maintain and consolidate these changes. The take-home exam continues to require a walkthrough video in which students reflect on the meaning of their code and demonstrate debugging skills. In addition, a final in-class presentation of the project has been formally integrated into the course structure. This presentation strengthens students’ ability to reflect on methodological and coding choices and fosters peer feedback. Both measures are intended to ensure that testing standards remain high, with a stronger emphasis on oral examination and critical reflection rather than only on basic coding skills. This focus has become increasingly important with the advent of AI-assisted coding tools, as it safeguards the integrity of assessment while keeping learning outcomes aligned with the course objectives.

\bibliographystyle{apacite}
\bibliography{../../resources/literature}

 
 
\end{document}
